% ============================================================
% Ejercicio 2
% ============================================================
\section{Ejercicio 2}

\begin{tcolorbox}[enunciado]
Proponer un certificado, junto con su esquema de codificación.
\begin{itemize}
    \item Describr e implemetar un algoritmo para generar certificados aleatorios para ejemplares. El programa debe de recibir como entrada (desde la línea de comandos):
    \begin{enumerate}
        \item nombre del archivo de entrada con la información del ejemplar.
        \item nombre del archivo de salida en donde se guardará el certificado generado.
    \end{enumerate}
\end{itemize}
\end{tcolorbox}

\subsection{Descripción del certificado y Algoritmo}
La codificación sigue la misma empleada en la práctica anterior, la codificación binaria. El algoritmo se presenta en el archivo: Generador Certificados.py. Quita los 32 bits del final de la cadena que corresponden al valor de W y de V, Para ir tomando las tuplas (wi,vi) y generar una cadena binaria que representa qué pedidos han sido tomados (1) y cuales no (0), por lo que esa cadena tendrá la misma cardinalidad que la cantidad de pedidos del ejemplar.
Toma un pedido con una probabilidad del 25%.